\chapter{基线、定位方法和取舍标准}

我们把“能不能跑完”和“跑得快不快”分开记录。前者使用功能测例和完整阶段回归确认，后者使用同一宿主、同一 QEMU、同一镜像和同一输入做交错 A/B。完整 BuildStorm 的计时从 workload 启动后的统一标记开始，到内核报告 \code{BUILDSTORM\_COMPILE mode=multi ok=true elapsed\_s} 为止；短 pc-hot 只用来筛选热点，不作为最终验收结果。

正式实验使用 QEMU 9.2.1、RISC-V QEMU virt、8 个 vCPU、16 GiB 内存和带 \code{-snapshot} 的发布镜像。每轮记录提交号、内核和镜像哈希、QEMU 参数、开始和结束时间以及完整日志路径。涉及架构无关代码时，我们至少执行 RISC-V 和 LoongArch 的静态检查；涉及文件系统、MM、task 生命周期的改动，还要执行对应的 final kernel 和功能回归。

我们采用一个很实际的取舍标准：完整 workload 有稳定收益且没有 panic、SIGSEGV、OOM、文件系统错误或测例缺项，才允许合入 \code{main}。低于宿主噪声的变化记为“无可确认收益”。例如，某些路径规范化和时间换算实验确实降低了局部指令数，但完整轮没有改善，所以没有放进最终方案。

\section{从现象到根因}
我们的定位过程大致固定为四步。第一步是在完整日志中确认首个异常或耗时阶段；第二步用计时和低开销计数确认调用次数、锁等待或页表操作是否真的集中；第三步做一个只改变单一因素的 A/B；第四步回到完整轮，检查局部收益是否能传递到端到端结果。这样做有点慢，却能避免把“热点”误当成“根因”。

在这条过程中，我们逐渐确认了几类重复成本：同一个路径在一次系统调用链中被多次解析，同一个文件的 metadata 被不同层重复读取；ELF 和 mmap 路径在没有实际 PTE 变化时仍进行同步；前台页分配承担了整页清零；而一些轻量查询仍然绕过同一份全局 registry 锁。后续章节按这些根因展开。
